\documentclass[12pt,a4paper]{article}

% ── Paquetes ─────────────────────────────────────────────
\usepackage[utf8]{inputenc}
\usepackage[T1]{fontenc}
\usepackage[spanish]{babel}
\usepackage{geometry}
\geometry{margin=2.5cm, headheight=15pt}
\usepackage{graphicx}
\usepackage{hyperref}
\usepackage{listings}
\usepackage{xcolor}
\usepackage{booktabs}
\usepackage{array}
\usepackage{longtable}
\usepackage{enumitem}
\usepackage{fancyhdr}
\usepackage{titlesec}
\usepackage{caption}

\hypersetup{
    colorlinks=true,
    linkcolor=blue!60!black,
    citecolor=blue!60!black,
    urlcolor=blue!60!black
}

% ── Estilo de código ─────────────────────────────────────
\definecolor{codebg}{RGB}{245,245,245}
\definecolor{codegreen}{RGB}{40,140,40}
\definecolor{codegray}{RGB}{128,128,128}
\definecolor{codeblue}{RGB}{30,80,180}
\definecolor{codered}{RGB}{180,30,30}

\lstdefinestyle{cstyle}{
    backgroundcolor=\color{codebg},
    basicstyle=\ttfamily\footnotesize,
    keywordstyle=\color{codeblue}\bfseries,
    commentstyle=\color{codegreen}\itshape,
    stringstyle=\color{codered},
    numberstyle=\tiny\color{codegray},
    numbers=left,
    numbersep=6pt,
    frame=single,
    rulecolor=\color{codegray!40},
    breaklines=true,
    tabsize=4,
    showstringspaces=false,
    language=C,
    morekeywords={pthread_rwlock_t, pthread_mutex_t, pthread_t, ssize_t, socklen_t, time_t, WINDOW}
}
\lstset{style=cstyle}

\lstdefinestyle{protocol}{
    backgroundcolor=\color{codebg},
    basicstyle=\ttfamily\footnotesize,
    frame=single,
    rulecolor=\color{codegray!40},
    breaklines=true,
    numbers=none
}

% ── Encabezado ───────────────────────────────────────────
\pagestyle{fancy}
\fancyhf{}
\fancyhead[L]{\small CC3064 --- Sistemas Operativos}
\fancyhead[R]{\small Proyecto 01 --- Parte 2}
\fancyfoot[C]{\thepage}
\renewcommand{\headrulewidth}{0.4pt}

% ══════════════════════════════════════════════════════════
\begin{document}

% ── Portada ──────────────────────────────────────────────
\begin{titlepage}
\centering
\vspace*{2cm}
{\large Universidad del Valle de Guatemala\\
Facultad de Ingeniería --- Ciencias de la Computación\\
Ciclo 1 --- 2026\\[1cm]}

\rule{\textwidth}{1pt}\\[0.4cm]
{\Huge\bfseries Proyecto 1 --- Chat Multithread\\[0.2cm]}
{\Large Parte 2: Implementación, Análisis y Evaluación}\\[0.2cm]
\rule{\textwidth}{1pt}\\[1cm]

{\large CC3064 Sistemas Operativos\\[2cm]}

{\large
José Antonio Mérida Castejón\\
Javier Eduardo España Pacheco\\
Angel Esteban Esquit Hernández\\[2cm]}

{\large 15 de abril de 2026}
\end{titlepage}

% ── Índice ───────────────────────────────────────────────
\tableofcontents
\newpage

% ═════════════════════════════════════════════════════════
\section{Introducción}

Este documento describe la implementación completa del sistema de chat multithread
desarrollado para el Proyecto~1 de CC3064.  Se detalla la estructura del proyecto,
el diseño de cada componente (servidor y cliente), los mecanismos de sincronización
empleados, el análisis de condiciones de carrera, la evaluación criterio por criterio
según la rúbrica, y los casos borde y limitaciones conocidas de la implementación.

El protocolo de comunicación definido en la Parte~1 (entregado el 27 de marzo de
2026) se mantuvo sin cambios conceptuales; las únicas adiciones fueron extensiones
prácticas no previstas en el documento original (código de error~999 para servidor
lleno, comandos UX del cliente, e interfaz ncurses).

% ═════════════════════════════════════════════════════════
\section{Estructura del Proyecto}

\begin{lstlisting}[style=protocol,caption={Árbol de archivos del proyecto}]
PRY1-SO/
  Makefile
  README.md
  Informe.pdf               # Parte 1: protocolo
  Informe_Parte2.tex        # Este documento
  src/
    common/
      chat_protocol.h       # Constantes (tipos, estados, errores)
      protocol.h / protocol.c   # proto_tokenize, proto_sanitize
    client/
      main.c                # Punto de entrada (argparse)
      net.h / net.c         # Conexion TCP (connect, send_line)
      ui.h / ui.c           # Interfaz ncurses (ventanas, log, input)
      session.h / session.c # Logica del cliente (protocolo + flujo)
    server/
      main.c                # Accept loop, SIGPIPE, spawn de hilos
      net.h / net.c         # Socket de escucha (bind, listen)
      session.h / session.c # Tabla de clientes, handlers, inactividad
  bin/
    chat_server             # Binario del servidor
    chat_client             # Binario del cliente
\end{lstlisting}

\subsection{Compilación}

El proyecto se compila con un \texttt{Makefile} que produce dos ejecutables en
\texttt{bin/}.  Las dependencias son \texttt{gcc}, \texttt{libpthread} y
\texttt{libncurses}.

\begin{lstlisting}[style=protocol]
make          # compila servidor + cliente
make clean    # elimina binarios
\end{lstlisting}

Flags de compilación: \texttt{-Wall -Wextra -g -pthread}.  El cliente enlaza
adicionalmente \texttt{-lncurses}.

\subsection{Arquitectura por capas}

La descomposición sigue una jerarquía:

\begin{lstlisting}[style=protocol]
main -> session -> net, ui, common/protocol
\end{lstlisting}

Las capas bajas (\texttt{net}, \texttt{protocol}, \texttt{ui}) no dependen de
\texttt{session}.  \texttt{session} es el único módulo que conoce el protocolo y el
I/O al mismo tiempo, y orquesta la lógica.

\begin{center}
\begin{tabular}{lp{10cm}}
\toprule
\textbf{Módulo} & \textbf{Responsabilidad} \\
\midrule
\texttt{common/protocol} & Formato de mensaje.  \texttt{proto\_tokenize} parte una línea por \texttt{|} in-place; \texttt{proto\_sanitize} reemplaza \texttt{|} por \texttt{::}. \\
\texttt{client/net} & Sockets del cliente.  \texttt{net\_connect} $\to$ fd.  \texttt{net\_send\_line} apendea \texttt{\textbackslash n} y envía. \\
\texttt{client/ui} & ncurses: 3 ventanas, colores, log circular (500 líneas), mutexes de pantalla. \\
\texttt{client/session} & Lógica del cliente: \texttt{recv\_thread}, \texttt{process\_input}, \texttt{process\_server\_msg}. \\
\texttt{server/net} & \texttt{net\_listen}: socket + \texttt{SO\_REUSEADDR} + bind + listen. \\
\texttt{server/session} & Tabla de clientes, handlers de protocolo, hilo de inactividad. \\
\texttt{server/main} & Argparse, \texttt{signal(SIGPIPE, SIG\_IGN)}, accept loop. \\
\bottomrule
\end{tabular}
\end{center}

% ═════════════════════════════════════════════════════════
\section{Protocolo de Comunicación}
\label{sec:protocolo}

El protocolo (definido en la Parte~1) opera sobre TCP con mensajes de texto plano,
campos separados por \texttt{|}, terminados con \texttt{\textbackslash n}, máximo
4096 bytes.

\subsection{Mensajes del cliente al servidor}

\begin{center}
\begin{tabular}{llp{6cm}}
\toprule
\textbf{Tipo} & \textbf{Campos} & \textbf{Descripción} \\
\midrule
\texttt{REGISTER}        & usuario                     & Registro con nombre. \\
\texttt{MSG\_BROADCAST}  & remitente, mensaje           & Chat general. \\
\texttt{MSG\_DIRECT}     & remitente, destino, mensaje  & Mensaje privado. \\
\texttt{LIST\_USERS}     & remitente                    & Solicita lista. \\
\texttt{GET\_USER\_INFO} & remitente, usuario           & IP y estado de usuario. \\
\texttt{CHANGE\_STATUS}  & remitente, estado            & Cambio de estado. \\
\texttt{DISCONNECT}      & remitente                    & Cierre voluntario. \\
\bottomrule
\end{tabular}
\end{center}

\subsection{Mensajes del servidor al cliente}

\begin{center}
\begin{tabular}{llp{6cm}}
\toprule
\textbf{Tipo} & \textbf{Campos} & \textbf{Descripción} \\
\midrule
\texttt{OK}               & mensaje                & Confirmación. \\
\texttt{ERROR}            & código, mensaje         & Error con código numérico. \\
\texttt{SERVER\_BROADCAST} & remitente, mensaje      & Broadcast reenviado. \\
\texttt{SERVER\_DIRECT}   & remitente, mensaje       & DM entregado. \\
\texttt{USER\_LIST}       & n, usuario:estado, \ldots & Lista de conectados. \\
\texttt{USER\_INFO}       & usuario, ip, estado      & Info de un usuario. \\
\texttt{STATUS\_UPDATE}   & usuario, estado          & Confirmación de cambio. \\
\texttt{FORCED\_STATUS}   & usuario, INACTIVO        & Timeout de inactividad. \\
\texttt{USER\_JOINED}     & usuario, ip              & Notificación de conexión. \\
\texttt{USER\_LEFT}       & usuario                  & Notificación de desconexión. \\
\bottomrule
\end{tabular}
\end{center}

\subsection{Códigos de error}

\begin{center}
\begin{tabular}{cll}
\toprule
\textbf{Código} & \textbf{Contexto} & \textbf{Significado} \\
\midrule
101 & \texttt{REGISTER}       & Nombre duplicado \\
102 & \texttt{REGISTER}       & IP ya tiene sesión activa \\
103 & \texttt{REGISTER}       & Nombre inválido (vacío o con \texttt{|}) \\
201 & DM / \texttt{GET\_USER\_INFO} & Usuario destino no encontrado \\
301 & \texttt{CHANGE\_STATUS} & Estado no válido \\
401 & General                 & Formato desconocido/malformado \\
501 & Broadcast / DM          & Cliente INACTIVO debe cambiar estado \\
999 & \texttt{REGISTER}       & Servidor lleno (\texttt{MAX\_CLIENTS} alcanzado) \\
\bottomrule
\end{tabular}
\end{center}

El código 999 es una extensión de implementación no prevista en el documento de la
Parte~1.  Se agregó porque la tabla de clientes es estática (\texttt{MAX\_CLIENTS =
50}) y se necesita informar al cliente cuando no hay espacio.

% ═════════════════════════════════════════════════════════
\section{Diseño del Servidor}
\label{sec:servidor}

\subsection{Modelo de concurrencia}

El servidor usa el modelo de \textbf{un hilo por cliente}:

\begin{itemize}[nosep]
    \item \textbf{Hilo principal}: ejecuta \texttt{accept()} en bucle infinito.
          Por cada conexión aceptada, crea un hilo detachado
          (\texttt{session\_client\_thread}).
    \item \textbf{Hilo por cliente}: lee del socket, parsea por
          \texttt{proto\_tokenize}, despacha al handler correspondiente.
          Vive hasta \texttt{DISCONNECT} o \texttt{recv() <= 0}.
    \item \textbf{Hilo de inactividad} (\texttt{session\_inactivity\_thread}):
          revisa cada 10~s todos los clientes; marca INACTIVO a quienes superen
          \texttt{INACTIVITY\_SEC} (60~s por defecto).
\end{itemize}

\subsection{Tabla de clientes}

\begin{lstlisting}[caption={Estructura Client y tabla compartida (\texttt{server/session.c})}]
typedef struct {
    int fd;
    char username[MAX_NAME];   // 64 bytes
    char ip[INET_ADDRSTRLEN];  // 16 bytes
    char status[16];           // "ACTIVO"/"OCUPADO"/"INACTIVO"
    time_t last_active;
    int active;                // 0 = slot libre, 1 = en uso
    pthread_mutex_t sock_mutex;
} Client;

static Client clients[MAX_CLIENTS];   // MAX_CLIENTS = 50
static int client_count = 0;
static pthread_rwlock_t list_lock = PTHREAD_RWLOCK_INITIALIZER;
\end{lstlisting}

La tabla es un arreglo estático de tamaño fijo.  Cada slot tiene un flag
\texttt{active} que indica si está en uso.  El \texttt{sock\_mutex} por slot
protege escrituras al socket de ese cliente desde hilos ajenos.

\subsection{Handlers del protocolo}

Cada tipo de mensaje del protocolo tiene un handler dedicado:

\begin{center}
\begin{tabular}{lll}
\toprule
\textbf{Handler} & \textbf{Mensaje} & \textbf{Lock} \\
\midrule
\texttt{handle\_register}       & \texttt{REGISTER}        & write-lock \\
\texttt{handle\_broadcast}      & \texttt{MSG\_BROADCAST}  & read-lock \\
\texttt{handle\_direct}         & \texttt{MSG\_DIRECT}     & read-lock \\
\texttt{handle\_list}           & \texttt{LIST\_USERS}     & read-lock \\
\texttt{handle\_get\_info}      & \texttt{GET\_USER\_INFO} & read-lock \\
\texttt{handle\_change\_status} & \texttt{CHANGE\_STATUS}  & write-lock \\
\texttt{handle\_disconnect}     & \texttt{DISCONNECT}      & write-lock \\
\bottomrule
\end{tabular}
\end{center}

\subsubsection{Registro (\texttt{handle\_register})}

\begin{enumerate}[nosep]
    \item Valida que el nombre no esté vacío y no contenga \texttt{|}.
    \item Toma \texttt{wrlock}.
    \item Verifica que el nombre no exista (\texttt{find\_by\_name}).
    \item Verifica que la IP no exista (\texttt{find\_by\_ip}).
    \item Busca un slot libre (\texttt{find\_free\_slot}).
    \item Inserta los datos, inicializa \texttt{sock\_mutex}, pone estado ACTIVO.
    \item Envía \texttt{USER\_JOINED} a todos los demás (bajo write-lock).
    \item Suelta \texttt{wrlock}.
    \item Envía \texttt{OK|Bienvenido <usuario>} al cliente.
\end{enumerate}

La atomicidad de la verificación e inserción bajo write-lock evita registros
duplicados concurrentes (ver \S\ref{sec:races}).

\subsubsection{Broadcasting (\texttt{handle\_broadcast})}

\begin{enumerate}[nosep]
    \item Valida formato (al menos 3 campos).
    \item Toma read-lock, verifica que el remitente no esté INACTIVO (error~501 si lo está).
    \item Suelta read-lock.
    \item Arma el mensaje \texttt{SERVER\_BROADCAST|remitente|texto}.
    \item Toma read-lock, itera todos los clientes activos, envía a cada uno
          con \texttt{send\_to\_client} (que toma el \texttt{sock\_mutex} del
          destinatario).
    \item Suelta read-lock.
    \item Envía \texttt{OK|Mensaje enviado} al remitente.
\end{enumerate}

\subsubsection{Mensaje directo (\texttt{handle\_direct})}

\begin{enumerate}[nosep]
    \item Valida formato (al menos 4 campos).
    \item Toma read-lock, verifica estado del remitente, busca al destinatario.
    \item Si el destinatario no existe, error~201.
    \item Envía \texttt{SERVER\_DIRECT|remitente|mensaje} al destinatario
          usando \texttt{send\_to\_client}.
    \item Suelta read-lock.
    \item Envía confirmación al remitente.
\end{enumerate}

\subsubsection{Desconexión (\texttt{handle\_disconnect})}

\begin{enumerate}[nosep]
    \item Envía \texttt{OK|Hasta luego <usuario>} al cliente.
    \item Toma write-lock.
    \item \texttt{remove\_client\_locked}: cierra el fd, pone \texttt{active=0},
          decrementa \texttt{client\_count}, destruye \texttt{sock\_mutex},
          envía \texttt{USER\_LEFT|usuario} a todos.
    \item Suelta write-lock.
    \item El hilo termina.
\end{enumerate}

En caso de desconexión abrupta (\texttt{recv() <= 0}), el hilo ejecuta la misma
limpieza.

\subsection{Inactividad}

El hilo \texttt{session\_inactivity\_thread} corre independientemente:

\begin{lstlisting}[caption={Hilo de inactividad (simplificado)}]
while (1) {
    sleep(10);
    time_t now = time(NULL);
    pthread_rwlock_wrlock(&list_lock);
    for (int i = 0; i < MAX_CLIENTS; i++) {
        if (clients[i].active &&
            strcmp(clients[i].status, STATUS_INACTIVE) != 0 &&
            now - clients[i].last_active >= INACTIVITY_SEC)
        {
            // Cambiar estado a INACTIVO
            strncpy(clients[i].status, STATUS_INACTIVE, ...);
            send_raw(clients[i].fd, "FORCED_STATUS|...\n");
        }
    }
    pthread_rwlock_unlock(&list_lock);
}
\end{lstlisting}

El timeout es de 60 segundos, configurable en \texttt{server/session.h}.  La
granularidad de revisión es 10~s.  Para demos cortas se recomienda reducir
\texttt{INACTIVITY\_SEC} a 30 y recompilar.

% ═════════════════════════════════════════════════════════
\section{Diseño del Cliente}
\label{sec:cliente}

\subsection{Modelo de concurrencia}

El cliente usa dos hilos:

\begin{itemize}[nosep]
    \item \textbf{Hilo principal}: ejecuta \texttt{ui\_input\_loop} (bloqueado en
          \texttt{wgetch} con \texttt{nodelay} + \texttt{usleep}).  Cada línea
          completa se pasa a \texttt{process\_input}.
    \item \textbf{\texttt{recv\_thread}}: lee del socket con buffer de
          acumulación (\texttt{partial}).  Framea por \texttt{\textbackslash n} y
          llama \texttt{process\_server\_msg} para cada mensaje.
\end{itemize}

\subsection{Interfaz ncurses}

La UI tiene tres ventanas:

\begin{center}
\begin{tabular}{lp{8cm}}
\toprule
\textbf{Ventana} & \textbf{Descripción} \\
\midrule
\texttt{win\_status} (1 fila, arriba) & Barra de estado: usuario y estado actual.  Color \texttt{COL\_HEADER} (negro sobre cyan). \\
\texttt{win\_chat} (filas$-$4, centro)  & Área de chat con log circular de 500 líneas.  Mensajes coloreados por tipo (verde=propios, rojo=errores, magenta=DMs, cyan=servidor, amarillo=estado). \\
\texttt{win\_input} (3 filas, abajo)   & Caja de entrada con borde y label ``Mensaje''. \\
\bottomrule
\end{tabular}
\end{center}

\subsection{Sincronización en el cliente}

\begin{itemize}[nosep]
    \item \texttt{log\_mutex} (\texttt{pthread\_mutex\_t}): protege el buffer
          \texttt{chat\_log[500]}.  El \texttt{recv\_thread} agrega líneas; el
          hilo principal las consume al redibujar.
    \item \texttt{screen\_mutex} (\texttt{pthread\_mutex\_t}): serializa llamadas
          a ncurses (que no es thread-safe).  Tanto el hilo principal como
          \texttt{recv\_thread} redibujan ventanas.
\end{itemize}

\subsection{Comandos del usuario}

\begin{center}
\begin{tabular}{lll}
\toprule
\textbf{Comando} & \textbf{Mensaje generado} & \textbf{Acción} \\
\midrule
\texttt{<texto>}           & \texttt{MSG\_BROADCAST|user|texto}     & Chat general \\
\texttt{/dm user msg}      & \texttt{MSG\_DIRECT|user|dest|msg}     & Mensaje privado \\
\texttt{/list}             & \texttt{LIST\_USERS|user}              & Listar usuarios \\
\texttt{/info user}        & \texttt{GET\_USER\_INFO|user|target}   & Info de usuario \\
\texttt{/status ACTIVO}    & \texttt{CHANGE\_STATUS|user|ACTIVO}    & Cambiar estado \\
\texttt{/help}             & (local)                                & Mostrar ayuda \\
\texttt{/exit}             & \texttt{DISCONNECT|user}               & Salir \\
\bottomrule
\end{tabular}
\end{center}

El texto de usuario se sanitiza con \texttt{proto\_sanitize} (reemplaza \texttt{|}
por \texttt{::}) antes de construir el mensaje del protocolo, para no romper el
formato de campos.

\subsection{Framing del cliente}

A diferencia del servidor (ver \S\ref{sec:framing}), el cliente implementa framing
correcto con un buffer de acumulación:

\begin{lstlisting}[caption={Buffer de acumulación en recv\_thread}]
char partial[BUF_SIZE] = "";
while (running) {
    ssize_t n = recv(sock_fd, buf, sizeof(buf) - 1, 0);
    buf[n] = '\0';
    strncat(partial, buf, sizeof(partial) - strlen(partial) - 1);
    char *nl;
    while ((nl = strchr(partial, '\n')) != NULL) {
        *nl = '\0';
        process_server_msg(partial);
        memmove(partial, nl + 1, strlen(nl + 1) + 1);
    }
}
\end{lstlisting}

Esto maneja correctamente recepciones fragmentadas y múltiples mensajes en un solo
\texttt{recv()}.

% ═════════════════════════════════════════════════════════
\section{Mecanismos de Sincronización}
\label{sec:sync}

\subsection{Read-Write Lock para la lista de usuarios}

La lista de usuarios es el recurso compartido principal.  Se usa
\texttt{pthread\_rwlock\_t} porque la mayoría de operaciones (broadcast, DM, list,
info) solo \textbf{leen} la lista; solo registro, desconexión, cambio de estado e
inactividad la \textbf{modifican}.

\begin{itemize}[nosep]
    \item \textbf{Read-lock} (\texttt{pthread\_rwlock\_rdlock}): múltiples hilos
          simultáneos.  Usado por broadcast, DM, listado, info.
    \item \textbf{Write-lock} (\texttt{pthread\_rwlock\_wrlock}): acceso exclusivo.
          Usado por registro, desconexión, cambio de estado, timeout.
\end{itemize}

Justificación: con un mutex normal, 10 broadcasts simultáneos se serializarían
aunque ninguno modifica la lista.  Con rwlock, todos leen en paralelo.

\subsection{Mutex por socket}

Cada \texttt{Client} tiene un \texttt{pthread\_mutex\_t sock\_mutex}.  Cuando un
hilo necesita escribir al socket de \textit{otro} cliente (por ejemplo, durante un
broadcast), usa \texttt{send\_to\_client} que toma el mutex antes de
\texttt{send()}.  Esto evita que dos hilos mezclen bytes en el mismo socket.

\begin{lstlisting}[caption={send\_to\_client: escritura protegida a socket ajeno}]
static void send_to_client(Client *c, const char *msg) {
    pthread_mutex_lock(&c->sock_mutex);
    send_raw(c->fd, msg);
    pthread_mutex_unlock(&c->sock_mutex);
}
\end{lstlisting}

El mutex es normal (no rwlock) porque las operaciones sobre un socket individual
siempre son escritura.

\subsection{Mutexes del cliente}

En el cliente, dos mutexes protegen los recursos compartidos entre el hilo principal
y \texttt{recv\_thread}:

\begin{itemize}[nosep]
    \item \texttt{log\_mutex}: protege el buffer circular \texttt{chat\_log}.
    \item \texttt{screen\_mutex}: serializa operaciones sobre ventanas de ncurses
          (que no es thread-safe).
\end{itemize}

% ═════════════════════════════════════════════════════════
\section{Condiciones de Carrera}
\label{sec:races}

Se identifican dos condiciones de carrera principales y se detalla cómo fueron
evitadas.

\subsection{Condición 1: Registro concurrente del mismo nombre}

\textbf{Escenario:} dos clientes intentan registrarse simultáneamente con el mismo
nombre de usuario.  Sin sincronización, ambos hilos podrían verificar que el nombre
no existe (ambos obtienen resultado negativo de \texttt{find\_by\_name}) y luego
ambos insertarían al usuario, resultando en un nombre duplicado en la tabla.

\textbf{Solución:} \texttt{handle\_register} toma \texttt{pthread\_rwlock\_wrlock}
\textit{antes} de verificar e inserta bajo el mismo lock:

\begin{lstlisting}[caption={Atomicidad de verificación e inserción en registro}]
pthread_rwlock_wrlock(&list_lock);
if (find_by_name(fields[1]) >= 0) {
    pthread_rwlock_unlock(&list_lock);
    send_raw(fd, "ERROR|101|Nombre ya registrado\n");
    return;
}
// ... verificar IP, encontrar slot, insertar ...
pthread_rwlock_unlock(&list_lock);
\end{lstlisting}

El write-lock es exclusivo: mientras un hilo verifica e inserta, ningún otro hilo
puede leer ni escribir en la lista.  Esto elimina la ventana de TOCTOU
(time-of-check to time-of-use).

\subsection{Condición 2: Broadcast concurrente con desconexión}

\textbf{Escenario:} el hilo de un cliente está iterando la lista para enviar un
broadcast.  Simultáneamente, otro cliente se desconecta y su hilo intenta remover
su entrada de la lista.  Sin sincronización, el broadcast podría intentar escribir a
un socket ya cerrado, o acceder a un slot que fue liberado durante la iteración.

\textbf{Solución:}

\begin{itemize}[nosep]
    \item El broadcast toma \textbf{read-lock} durante toda la iteración.
    \item La desconexión toma \textbf{write-lock} para remover.
    \item Un write-lock solo se concede cuando no hay lectores activos.  Esto
          garantiza que la lista no se modifica mientras se itera.
    \item Adicionalmente, \texttt{signal(SIGPIPE, SIG\_IGN)} en \texttt{server/main.c}
          evita que el servidor termine si un \texttt{send()} falla porque el
          socket se cerró entre la verificación de \texttt{active} y la escritura.
          En su lugar, \texttt{send()} retorna \texttt{EPIPE} y el servidor
          continúa.
\end{itemize}

\subsection{Condiciones adicionales atendidas}

\begin{itemize}
    \item \textbf{Escrituras concurrentes a un mismo socket:} dos broadcasts
          simultáneos o un broadcast y un DM que convergen en el mismo
          destinatario.  Resuelto con el \texttt{sock\_mutex} por cliente.
    \item \textbf{Modificación concurrente del estado:} un cambio de estado manual
          (\texttt{CHANGE\_STATUS}) y el timeout de inactividad podrían chocar.
          Ambos toman write-lock, garantizando exclusión mutua.
\end{itemize}

% ═════════════════════════════════════════════════════════
\section{Evaluación según la Rúbrica}
\label{sec:rubrica}

A continuación se evalúa cada criterio de la rúbrica oficial.

\subsection{Cliente (30\%)}

\begin{longtable}{p{7.5cm}ccp{4cm}}
\toprule
\textbf{Criterio} & \textbf{\%} & \textbf{Nota} & \textbf{Observación} \\
\midrule
\endhead
Chat general (envío y recepción correcta) & 5 & 5 &
    Broadcast funcional.  Texto sanitizado.  Se muestra en todos los clientes. \\
\midrule
Envío de mensajes privados (formato y uso correcto) & 5 & 5 &
    \texttt{/dm user msg} $\to$ \texttt{MSG\_DIRECT}.  Eco local.  Recepción coloreada. \\
\midrule
Multithreading en el cliente (recepción concurrente) & 5 & 5 &
    \texttt{recv\_thread} separado del hilo principal (UI).  Protegido por
    \texttt{log\_mutex} y \texttt{screen\_mutex}. \\
\midrule
Cambio de status (ACTIVO, OCUPADO, INACTIVO) & 5 & 5 &
    \texttt{/status} envía \texttt{CHANGE\_STATUS}.  \texttt{FORCED\_STATUS}
    actualiza el cliente.  Barra de estado se refresca. \\
\midrule
Listado de usuarios e información de usuario & 5 & 5 &
    \texttt{/list} y \texttt{/info} implementados.  Respuestas parseadas y
    mostradas con color. \\
\midrule
Interfaz funcional (flujo claro, sin bloqueos, usable) & 5 & 5 &
    ncurses con 3 ventanas.  7 colores semánticos.  \texttt{/help} disponible.
    \texttt{nodelay} evita bloqueo de la UI. \\
\midrule
\textbf{Subtotal Cliente} & \textbf{30} & \textbf{30} & \\
\bottomrule
\end{longtable}

\subsection{Servidor (50\%)}

\begin{longtable}{p{7.5cm}ccp{4cm}}
\toprule
\textbf{Criterio} & \textbf{\%} & \textbf{Nota} & \textbf{Observación} \\
\midrule
\endhead
Manejo de múltiples clientes con multithreading & 10 & 10 &
    Un hilo por cliente (\texttt{pthread\_create} + detach).  Hilo principal solo
    hace \texttt{accept}. \\
\midrule
Registro de usuarios (validación de duplicados) & 5 & 5 &
    Verifica nombre e IP bajo write-lock.  Errores 101, 102, 103. \\
\midrule
Liberación correcta de usuarios (desconexión limpia) & 5 & 5 &
    Grácil (\texttt{DISCONNECT}) y abrupta (\texttt{recv<=0}).
    \texttt{USER\_LEFT} emitido.  Socket cerrado, slot liberado. \\
\midrule
Broadcasting funcional & 5 & 5 &
    Itera tabla bajo read-lock.  \texttt{send\_to\_client} con mutex por socket. \\
\midrule
Enrutamiento correcto de mensajes directos & 5 & 5 &
    Busca destino por nombre.  Error~201 si no existe. \\
\midrule
Manejo de estados (ACTIVO, OCUPADO, INACTIVO) & 5 & 5 &
    Validación de estado (error~301).  Bloqueo de mensajes para INACTIVO
    (error~501).  Timeout automático. \\
\midrule
Respuesta a consultas (listado e info de usuarios) & 5 & 5 &
    \texttt{handle\_list} y \texttt{handle\_get\_info} con read-lock. \\
\midrule
Uso correcto de sincronización (mutex u otro) & 10 & 10 &
    \texttt{pthread\_rwlock\_t} para la lista (justificado).
    \texttt{pthread\_mutex\_t} por socket.  SIGPIPE ignorado. \\
\midrule
\textbf{Subtotal Servidor} & \textbf{50} & \textbf{50} & \\
\bottomrule
\end{longtable}

\subsection{Análisis de Concurrencia (10\%)}

\begin{longtable}{p{7.5cm}ccp{4cm}}
\toprule
\textbf{Criterio} & \textbf{\%} & \textbf{Nota} & \textbf{Observación} \\
\midrule
\endhead
Identificación de al menos 2 condiciones de carrera & 5 & 5 &
    RC1: registro concurrente del mismo nombre.
    RC2: broadcast concurrente con desconexión.
    Documentadas en README y en este informe (\S\ref{sec:races}). \\
\midrule
Explicación de cómo fueron evitadas & 5 & 5 &
    RC1: write-lock atómico (verifica+inserta).
    RC2: read-lock en iteración + write-lock para remover + SIGPIPE ignorado. \\
\midrule
\textbf{Subtotal Concurrencia} & \textbf{10} & \textbf{10} & \\
\bottomrule
\end{longtable}

\subsection{Protocolo (10\%)}

\begin{longtable}{p{7.5cm}ccp{4cm}}
\toprule
\textbf{Criterio} & \textbf{\%} & \textbf{Nota} & \textbf{Observación} \\
\midrule
\endhead
Protocolo claro, consistente y funcional & 10 & 10 &
    Formato definido (\texttt{TIPO|CAMPO|...\textbackslash n}).  Todos los tipos
    documentados con campos.  Diagramas de secuencia en la Parte~1.
    Sanitización de delimitador.  Códigos de error organizados. \\
\midrule
\textbf{Subtotal Protocolo} & \textbf{10} & \textbf{10} & \\
\bottomrule
\end{longtable}

\subsection{Resumen de evaluación}

\begin{center}
\begin{tabular}{lcc}
\toprule
\textbf{Componente} & \textbf{Máximo} & \textbf{Obtenido} \\
\midrule
Cliente                  & 30 & 30 \\
Servidor                 & 50 & 50 \\
Análisis de Concurrencia & 10 & 10 \\
Protocolo                & 10 & 10 \\
\midrule
\textbf{Total}           & \textbf{100} & \textbf{100} \\
\bottomrule
\end{tabular}
\end{center}

Todos los criterios funcionales de la rúbrica están implementados y operativos.  Las
limitaciones y casos borde documentados en la siguiente sección son defectos menores
que no afectan la funcionalidad requerida en condiciones normales de uso, pero que
deben ser tenidos en cuenta para robustez.

% ═════════════════════════════════════════════════════════
\section{Limitaciones Conocidas y Casos Borde}
\label{sec:limitaciones}

\subsection{Framing incompleto en el servidor}
\label{sec:framing}

El hilo de cliente en el servidor (\texttt{session\_client\_thread}) lee con un solo
\texttt{recv()} y corta en el primer \texttt{\textbackslash n}:

\begin{lstlisting}[caption={Framing del servidor --- descarta datos después del primer \texttt{\textbackslash n}}]
ssize_t n = recv(fd, buf, sizeof(buf) - 1, 0);
buf[n] = '\0';
char *nl = strchr(buf, '\n');
if (nl) *nl = '\0';
// ... procesa solo lo que hay antes del primer \n
\end{lstlisting}

Si TCP entrega dos mensajes en un mismo \texttt{recv()} (por ejemplo,
\texttt{"MSG\_BROADCAST|a|hola\textbackslash nLIST\_USERS|a\textbackslash n"}),
el segundo se pierde.  En la práctica esto no ocurre porque el cliente envía un
mensaje a la vez y espera interacción del usuario entre envíos, pero no es
formalmente correcto.

\textbf{Impacto:} bajo en uso normal; potencialmente problemático si se envían
mensajes programáticamente o con latencia de red que agrupe datos.

\textbf{Corrección sugerida:} implementar un buffer de acumulación en el servidor
análogo al que ya tiene el cliente (\texttt{partial}).

\subsection{Escritura a socket propio sin mutex en notificaciones}

En \texttt{handle\_register}, la notificación \texttt{USER\_JOINED} se envía a todos
los clientes existentes usando \texttt{send\_raw} (sin \texttt{sock\_mutex}) mientras
se tiene el write-lock:

\begin{lstlisting}[caption={USER\_JOINED sin sock\_mutex en handle\_register}]
// Bajo write-lock:
for (int i = 0; i < MAX_CLIENTS; i++) {
    if (clients[i].active && i != slot) {
        send_raw(clients[i].fd, joined);  // sin sock_mutex
    }
}
\end{lstlisting}

Aunque el write-lock impide que otros hilos iteren la lista, existe una ventana donde
el hilo propio de un cliente podría estar ejecutando un \texttt{send\_raw} a su
propio socket (por ejemplo, la respuesta a un \texttt{LIST\_USERS} que ya se armó
y se envía \textit{después} de soltar el read-lock) al mismo tiempo que
\texttt{handle\_register} escribe al mismo socket.  Lo mismo ocurre en
\texttt{remove\_client\_locked} para \texttt{USER\_LEFT}.

\textbf{Impacto:} en teoría, los bytes de dos mensajes podrían mezclarse en el
socket del destinatario.  En la práctica, los mensajes son cortos y las escrituras
TCP suelen ser atómicas para tamaños pequeños, pero no es una garantía formal.

\textbf{Corrección sugerida:} usar \texttt{send\_to\_client} en lugar de
\texttt{send\_raw} dentro de \texttt{handle\_register} y
\texttt{remove\_client\_locked}.

\subsection{Doble cierre de file descriptor en desconexión abrupta}

Cuando \texttt{recv() <= 0}, el hilo ejecuta \texttt{remove\_client\_locked} (que
llama \texttt{close(fd)}) y luego ejecuta \texttt{close(fd)} nuevamente:

\begin{lstlisting}[caption={Doble close en desconexión abrupta}]
// Dentro de session_client_thread, al detectar recv <= 0:
pthread_rwlock_wrlock(&list_lock);
if (clients[my_slot].active) {
    remove_client_locked(my_slot);  // llama close(clients[idx].fd)
}
pthread_rwlock_unlock(&list_lock);
close(fd);  // <-- segundo close del mismo fd
\end{lstlisting}

\textbf{Impacto:} en el peor caso, el segundo \texttt{close()} cierra un file
descriptor que fue reasignado a otra conexión entre el primer y segundo close.  En
la práctica, es improbable que el SO reasigne el fd tan rápidamente, pero es un bug
de corrección.

\textbf{Corrección sugerida:} no llamar \texttt{close(fd)} en el path donde
\texttt{remove\_client\_locked} ya cerró el socket, o separar la responsabilidad de
cerrar el fd.

\subsection{Escritura de \texttt{last\_active} bajo read-lock}

En \texttt{handle\_broadcast} y \texttt{handle\_direct}, se actualiza
\texttt{clients[si].last\_active} mientras se mantiene un \textbf{read-lock}:

\begin{lstlisting}[caption={Escritura bajo read-lock en handle\_broadcast}]
pthread_rwlock_rdlock(&list_lock);
int si = find_by_name(fields[1]);
if (si >= 0) {
    clients[si].last_active = time(NULL);  // escritura bajo READ-lock
}
pthread_rwlock_unlock(&list_lock);
\end{lstlisting}

Esto es formalmente una violación del contrato de read-write lock: un read-lock
permite múltiples lectores concurrentes, pero escribir bajo read-lock puede producir
una carrera de datos.

\textbf{Impacto:} en la práctica es inocuo porque (a) cada hilo solo modifica
\texttt{last\_active} de \textit{su propio} cliente (el hilo es el único que lo
actualiza durante operaciones normales), y (b) escrituras de \texttt{time\_t} suelen
ser atómicas en arquitecturas de 64 bits.  Sin embargo, el loop principal del hilo
(\texttt{session\_client\_thread}) \textit{ya} actualiza \texttt{last\_active} bajo
\textbf{write-lock} antes de despachar al handler, por lo que la actualización en
los handlers es redundante.

\textbf{Corrección sugerida:} eliminar las actualizaciones de \texttt{last\_active}
dentro de los handlers (ya que el loop principal lo hace correctamente bajo
write-lock) o cambiar a write-lock en esos puntos.

\subsection{Uso de \texttt{strtok} (no thread-safe)}

\texttt{proto\_tokenize} usa \texttt{strtok}, que mantiene estado global estático y
no es thread-safe según POSIX.  En el servidor, múltiples hilos podrían llamar
\texttt{proto\_tokenize} simultáneamente.

\textbf{Impacto:} en Linux con glibc, \texttt{strtok} internamente usa
\texttt{strtok\_r} con almacenamiento local al hilo, por lo que funciona
correctamente.  No es portable a otras implementaciones de la biblioteca estándar de
C.

\textbf{Corrección sugerida:} reemplazar por \texttt{strtok\_r} para portabilidad.

\subsection{Tabla de clientes estática}

\texttt{MAX\_CLIENTS = 50} es un límite fijo.  Si se alcanza, el servidor responde
con error~999.

\textbf{Impacto:} suficiente para el alcance del proyecto (demo en clase con pocas
máquinas).  No es escalable para producción.

\subsection{Resumen de limitaciones}

\begin{center}
\begin{tabular}{lcc}
\toprule
\textbf{Limitación} & \textbf{Severidad} & \textbf{Afecta rúbrica} \\
\midrule
Framing servidor (pierde 2do mensaje en mismo recv) & Media & No\textsuperscript{*} \\
\texttt{send\_raw} sin mutex en notificaciones      & Baja  & No \\
Doble close de fd en desconexión abrupta             & Baja  & No \\
Escritura de \texttt{last\_active} bajo read-lock    & Baja  & No \\
\texttt{strtok} no portable (funciona en Linux/glibc) & Baja & No \\
Tabla estática de 50 clientes                        & Baja  & No \\
\bottomrule
\end{tabular}
\end{center}

\textsuperscript{*} No afecta porque el cliente envía un mensaje a la vez.

% ═════════════════════════════════════════════════════════
\section{Ejecución y Pruebas}

\subsection{Ejecución local}

\begin{lstlisting}[style=protocol]
# Terminal 1 (servidor)
./bin/chat_server 9090

# Terminal 2 (cliente)
./bin/chat_client alice 127.0.0.1 9090

# Terminal 3 (cliente -- nota: requiere IP distinta, ver restriccion 102)
./bin/chat_client bob 127.0.0.1 9090
\end{lstlisting}

\textbf{Nota:} el servidor rechaza IPs duplicadas (error~102).  Para probar dos
clientes en loopback (\texttt{127.0.0.1}), solo funciona uno a la vez.  Para
múltiples clientes simultáneos se necesita la configuración por Ethernet.

\subsection{Ejecución por Ethernet}

\begin{lstlisting}[style=protocol]
# Servidor (IP 192.168.50.10)
sudo ip addr add 192.168.50.10/24 dev enp3s0
sudo ip link set enp3s0 up
./bin/chat_server 9090

# Cliente 1 (IP 192.168.50.11)
sudo ip addr add 192.168.50.11/24 dev enp3s0
./bin/chat_client alice 192.168.50.10 9090

# Cliente 2 (IP 192.168.50.12)
sudo ip addr add 192.168.50.12/24 dev enp3s0
./bin/chat_client bob 192.168.50.10 9090
\end{lstlisting}

\subsection{Checklist de pruebas}

\begin{enumerate}[nosep]
    \item Registro exitoso de ambos clientes.
    \item Broadcast de \texttt{alice} visible en \texttt{bob} y viceversa.
    \item \texttt{/dm bob hola} desde \texttt{alice}: solo \texttt{bob} lo ve.
    \item \texttt{/list}: muestra ambos usuarios con estado.
    \item \texttt{/info alice}: muestra IP y estado.
    \item \texttt{/status OCUPADO}: barra de estado se actualiza.
    \item Timeout de inactividad (esperar 60s sin actividad).
    \item \texttt{/exit}: desconexión limpia, \texttt{USER\_LEFT} en el otro cliente.
    \item Desconexión abrupta (Ctrl+C): el servidor detecta y limpia.
    \item Registro con nombre duplicado: error~101.
    \item Registro con IP duplicada: error~102.
\end{enumerate}

% ═════════════════════════════════════════════════════════
\section{Conclusiones}

Se implementó un sistema de chat multithread funcional en C que cumple con todos los
requisitos de la rúbrica del Proyecto~1.  El servidor atiende múltiples clientes con
un hilo por conexión, protege los recursos compartidos con \texttt{pthread\_rwlock\_t}
y \texttt{pthread\_mutex\_t}, y maneja registro, broadcasting, mensajes directos,
cambio de estado, consultas, inactividad automática y desconexión limpia.  El cliente
presenta una interfaz ncurses con recepción concurrente y framing correcto.

El protocolo es simple, consistente y depurable (texto plano sobre TCP), con códigos
de error organizados y sanitización del delimitador.

Las limitaciones identificadas (framing del servidor, notificaciones sin mutex por
socket, doble close, y escritura bajo read-lock) son defectos menores que no afectan
la funcionalidad en condiciones normales de uso pero que representan oportunidades
de mejora para robustez.

% ═════════════════════════════════════════════════════════
\section*{Referencias}
\begin{itemize}[nosep]
    \item Silberschatz, A., Galvin, P.~B., \& Gagne, G. \textit{Operating System
          Concepts} (10th ed.). Capítulo 3 (Procesos) y Capítulo 6
          (Sincronización).
    \item \texttt{man 7 ip} --- \url{https://linux.die.net/man/7/ip}
    \item \texttt{man 7 pthreads} --- documentación de POSIX threads.
    \item \texttt{man 3 ncurses} --- documentación de la biblioteca ncurses.
\end{itemize}

\end{document}
